\documentclass{article}
\usepackage[utf8]{inputenc}
\begin{document}

\section{Exercise 3: Grammar}

This exercise is intended to teach strategies for identifying problems with standard grammar, punctuation and sentence-level errors. 

\subsection*{Part A}

\subsubsection*{Instructions}

\begin{itemize}
\item Read the following paragraphs.
\item There are 10 different mistakes present. These include issues with active/passive voice, tense conflicts, weak verbs, excessive use of prepositions, misused pronouns, improper word usage and incorrect punctuation. 
\item Highlight/bold/italicise these mistakes.
\item There's no need to fix the mistakes in this part - Part B will involve fixing grammatical errors in a different passage.
\item The answers are available on page 2 - but don't peek! 
\end{itemize}

Note: Discuss with your neighbor - science is a collaborative process! If you are running into any problems with Overleaf or have questions, please feel free to raise your hand and ask for help. 

\subsubsection*{Exercise 3: Part A}

Affect of different diets on the metabolic rate and respiratory quotient for \textit{Tenebrio molitor} 

\par Tenebrio molitor is also known as the darkling beetle; which exists in its larval stage as the yellow mealworm (Fraenkel, 1950). Similar to other living organisms, it carries out cellular respiration. This is a vital process, of physiology, and involves the oxidation of organic molecules, resulting in the release of the energy carrier ATP, carbon dioxide and water (Fenton et al., 2015). The following is an example of cellular respiration: 

C6H12O6 + 6 O2 ---------> 6 H2O + 6 CO2 + ATP 

However, there respiration is not limited to only carbohydrates. Other substrates that may also be respired by Tenebrio molitor include fats, organic acids and proteins. 

Wild Tenebrio molitor require diets with a high proportion of carbohydrates, cholesterol, and proteins such as casein and lactalbumin, for optimal growth (Fraenkel, 1950). 
 
The purpose of this lab was: to investigate the effect of the independent variable (the diet) on both the metabolic rate and the respiratory quotient of the Tenebrio molitor mealworms. The different diets used were water (‘starved’), a high protein diet and a regular lab diet of wheat, bran and oatmeal (i.e. rich in carbohydrates). 


ewpage
\subsubsection*{Part A - Solutions}

\begin{enumerate}
\item In the title, it should be effect - not affect.
\item Need to italicise the genus and species name (\textit{Tenebrio molitor})
\item Improper use of a semicolon in the first sentence. A conjunction is present (which), so a comma should have been used.
\item In the second sentence, the use of the word 'it' is ambiguous. Is the author referring to the yellow mealworm (larval stage) or the darkling beetle (adult form)? This is an example of a misused pronoun.
\item Too many prepositions are present in the third statement. E.g. "of physiology", "of the energy carrier"
\item Additionally, the third statement is a run-on sentence - i.e. there is an excessive use of commas.
\item Equation is typed incorrectly. Numbers should be subscript. (This isn't technically a grammatical mistake - but you should have spotted it!)
\item Should be "their", not there. (But if we're more picky, there's no need for a "there" in the sentence at all.)
\item "Other substrates ...." is an example of a passive voice sentence. There is also a weak verb present ("to be"). This sentence could be written much more concisely using the active voice.
\item No need for a colon in the sentence "The purpose of this lab was: to investigate."
\end{enumerate}


ewpage
\subsection*{Part B}

\subsubsection*{Instructions}

\begin{itemize}
\item This entire exercise revolves around sentence-level errors, caused by misplaced/excessive commas, colons and semi-colons. 
\item You will encounter problems such as sentence fragments, run-on sentences, pronoun disagreement and faulty parallelism. Don't worry if these sound unfamiliar - you have all the tools to edit these sentences and fix the underlying problems.
\item Solutions may involve having to split sentences, combine sentences or eliminate entire phrases. Make sure that you preserve the original meaning of the sentence.
\item Suggested solutions are available on page 4. 
\end{itemize}

Note: Discuss with your neighbor - science is a collaborative process! If you are running into any problems with Overleaf or have questions, please feel free to raise your hand and ask for help. 

\subsubsection*{Exercise 3: Part B}

\begin{enumerate}
\item Being able to successfully identify and distinguish microorganisms from each other is an important technique and is used in various important fields. For example, medical microbiology.
\item Generally, the first step of a traditional diagnostic method is to measure speed, colour and morphology.
\item The Gram’s stain procedure is a useful starting point; as the tests immediately narrows down the possibilities to either Gram positive or negative. 
\item The purpose of this experiment was to correctly determine the identity (i.e. genus and species) of the unknown bacterium sample, based on its structure, morphology, physiological and metabolic characteristics, using first-stage and second-stage diagnostic tests, tables and literature descriptions of bacterial genera. 
\item The acid-fast stain results in the dull bacterial cells stained blue. 
\end{enumerate}



ewpage 
\subsubsection*{Exercise 3: Solutions}

\begin{enumerate}
\item The problem here is "For example, medical microbiology." This is an example of a sentence fragment. The second sentence cannot stand alone - it would be better to combine it with the first.
\paragraph{} "Being able to successfully identify and distinguish microorganisms from each other is an important technique and is used in various important fields, such as medical microbiology."

\item This is an example of faulty parallelism. You cannot measure colour and morphology - both of these properties need a different verb (such as observe/record). If your sentence contains a list of ideas/items, with only one verb, make sure that the action is applicable to each idea. 
\paragraph{} "Generally, the first step of a traditional diagnostic method is to measure speed, observe colour and record the morphology."

\item A semi-colon should have not been used since the conjunction "as" was present. Additionally, \textit{test} should be singular in order to agree with the verb \textit{narrows}.
\paragraph{} "The Gram’s stain procedure is a useful starting point, as the test immediately narrows down the possibilities to either Gram positive or negative."

\item Trick question! There is no grammatical error in this one. But it is a run-on sentence - there are excessive clauses and commas. This sentence would be better off split into multiple shorter sentences.
\item This sentence is suffering from a tense conflict (should be \textit{resulted}) and a missing verb (needs "\textit{being}").
\paragraph{} "The acid-fast stain resulted in the dull bacterial cells being stained blue." 
\end{enumerate}



ewpage
\subsection*{Part C}

\subsubsection*{Instructions}

\begin{itemize}
\item Read the following paragraphs.
\item What mistakes can you spot? Edit as you go. If you find a way to be more concise, go ahead and do so!
\item Suggested solutions can be found on the next page (page 4) - but remember: there are multiple ways to fix these grammatical mistakes. 
\end{itemize}

Note: Discuss with your neighbor - science is a collaborative process! If you are running into any problems with Overleaf or have questions, please feel free to raise your hand and ask for help. 

\subsubsection*{Exercise 3: Part C}

The muscle is an important contractile organ system; that is composed of several proteinaceous components, which amongst other functions, it primarily serves to drive movement in an organism. Their are three types of muscle; smooth, cardiac, skeletal. 

Within a muscle, each muscle fibre contains packed myofibrils, which in turn, consist of sarcomeres, which are a functional contractile unit. Sarcomeres themselves contain adjacent myosin (thick) and actin (thin) filaments, resulting in a striated appearance [1]. The filaments are arranged as follows. The actin filaments are attached to a Z-disc and extend to interdigitate with myosin filaments [1, 2]. The Z-disc attaches myofibrils to each other, which aids in combined contraction [1]. Within muscle fibres is also a sarcolemma, nucleus, sarcoplasm and other proteins (such as titin, filamin and spectrin).

With the beginning of a cycle of contraction, myosin and actin are firmly interlocked. In the presence of ATP, it is hydrolysed by the myosin head and resulting in the two filaments being separated. The power stroke then occurs, e.g. a change in conformation, generating the force needed for myosin filaments slide along actin, causing the sarcomere to shorten [2]. The collective contraction of all the sarcomeres in a muscle fibre results in the entire muscle contracting as one unit. 


ewpage
\subsubsection*{Suggested Solution}

The muscle is an important contractile organ system\textbf{, which} is composed of several proteinaceous components\textbf{.} \textbf{The muscle }primarily serves to drive movement in an organism. \textbf{There} are three types of muscle\textbf{:} smooth, cardiac\textbf{ and }skeletal. 

\textbf{Muscles are composed of} muscle fibres, which \textbf{consist of }packed myofibrils. \textbf{These myofibrils }consist of sarcomeres \textbf{(a functional contractile unit)}. Sarcomeres \textbf{contain} adjacent myosin (thick) and actin (thin) filaments, resulting in a striated appearance [1]. The filaments are arranged \textbf{as follows: actin} filaments are attached to a Z-disc and extend to interdigitate with myosin filaments [1, 2]. The Z-disc attaches myofibrils to each other, which aids in combined contraction [1]. \textbf{Muscle fibres also contain} a sarcolemma, nucleus, sarcoplasm and other proteins (such as titin, filamin and spectrin).

\textbf{At} the beginning of \textbf{a contraction cycle}, myosin and actin are firmly interlocked. \textbf{If ATP is present, it will be hydrolysed by the myosin head, resulting in the separation of the two filaments}. The power stroke then occurs \textbf{(i.e. a change in conformation)} generating the force needed for myosin filaments slide along actin, causing the sarcomere to shorten [2]. \textbf{Collective sarcomere contraction} in a muscle fibre results in the entire muscle contracting as one unit. 

\end{document}